\chapter{平台控制、反馈与可靠性}

\begin{keyidea}
整机由大量不同时间尺度的控制环维持稳定。本章把硬件状态机、固件策略、驱动调度、反馈控制和故障恢复统一成平台控制面。
\end{keyidea}

前面几章把数据通路讲完了：最小 CPU 章给出了取指、译码、执行的五段路径，存储与互连相关章节给出了地址、数据、握手三类总线事务，存储器件相关章节给出了从 SRAM 单元到 DRAM 命令流的完整结构。但一台整机里还有第二张网：CPU 写出一个物理地址，磁盘控制器就开始寻道；设备拉出一条中断线，CPU 就放下当前指令转入服务程序；网卡收到一个包，不经任何一条 CPU 指令就出现在内存里。这些“什么时候让谁动”的决定，构成控制面。值得注意的是，CPU 与外界的一切交互只归结为三类信号：对物理地址的读与写（取指、访存、设备寄存器读写都是同一类事务）、从外部进入的中断输入、设备与内存之间的 DMA 搬运。本章把这张控制网讲透：先看任何设备共有的寄存器接口与状态机，再看 CPU 如何控制自己最快的从设备——内存，然后沿中断链路上行到固件与操作系统，最后把这些分散的机制统一成同一个概念——反馈环。

\begin{keyidea}
控制 = 状态机 + 反馈环。每个控制器（设备、内存控制器、中断控制器）都是一台状态机：它有一组可被软件读写的寄存器作为输入输出接口，有一组时序约束决定状态迁移何时合法。每个控制器又都处在至少一个反馈环里：发出动作、测量结果、按偏差调整。反馈环按响应时间分层——纳秒级的快环（DRAM 时序计数器、PHY 采样点、链路时钟恢复 CDR）必须由硬件闭合，因为软件根本来不及参与；微秒到秒级的慢环（驱动程序的设备管理、固件的内存训练、操作系统的调度）才交给软件。划分的原则只有一条：环的响应时间必须短于被控对象变化的时间常数。
\end{keyidea}

\section{一、控制面：数据通路之外的第二张网}

最小 CPU 章的数据通路图里，每一条线都有明确的方向和内容：PC 到指令存储器是地址，指令存储器到译码器是指令编码，ALU 到寄存器堆是结果。数据通路回答“数据怎么流动”，但它不回答三个更外层的问题：CPU 怎么让一台设备开始工作？设备怎么让 CPU 知道自己做完了？大批量数据怎么在设备与内存之间移动而不占用流水线？这三个问题的答案都不在数据通路里，而在控制面里。

控制面建立在同一个物理地址空间之上。存储与互连相关章节已经说明，MMIO（Memory-Mapped I/O，内存映射 I/O）把设备寄存器映射进物理地址空间，CPU 用普通的 load/store 就能访问设备；本节要做的是把视角反转过来：从设备一侧看，这些读写不是“访问内存”，而是命令与状态采样。三类信号各管一件事——对物理地址的读与写是一类：MMIO 写是 CPU 发出的命令，MMIO 读是 CPU 对设备状态的采样（取指与访存也属这类事务）；中断是设备反方向发来的通知；DMA 是数据绕开 CPU 流水线的搬运通道。

\begin{pdfdiagram}[x=1cm,y=1cm]
  % CPU 核心居中
  \node[hw compute, minimum width=2.4cm, minimum height=1.4cm] (cpu) at (5.5,0) {CPU 核心\\（流水线/Cache）};
  % 四颗卫星
  \node[hw memory, minimum width=2.6cm, minimum height=1.2cm] (mc) at (1.6,2.6) {内存控制器\\+ DRAM};
  \node[hw control, minimum width=2.6cm, minimum height=1.2cm] (apic) at (9.4,2.6) {中断控制器\\（LAPIC/IOAPIC）};
  \node[hw state, minimum width=2.6cm, minimum height=1.2cm] (pch) at (1.6,-2.6) {PCIe 根复合体\\/ PCH};
  \node[hw state, minimum width=2.6cm, minimum height=1.2cm] (dev) at (9.4,-2.6) {设备 / DMA\\（网卡/磁盘/UART）};
  % 内存读写: CPU -> MC
  \draw[hw bus] (4.3,0.35) -- (3.4,0.35) -- (3.4,2.6) -- (2.9,2.6);
  \node[hw label, font=\tiny, anchor=west, text width=2.4cm] at (3.55,1.5) {内存读写\\（物理地址）};
  % 中断: APIC -> CPU
  \draw[hw return] (8.1,2.6) -- (7.6,2.6) -- (7.6,0.35) -- (6.7,0.35);
  \node[hw label, font=\tiny, anchor=west] at (7.75,1.5) {中断请求};
  % MMIO: CPU -> PCH
  \draw[hw bus] (4.3,-0.35) -- (3.2,-0.35) -- (3.2,-2.6) -- (2.9,-2.6);
  \node[hw label, font=\tiny, anchor=west, text width=2.4cm] at (3.35,-1.5) {MMIO 读写\\（设备寄存器）};
  % 中断: 设备 -> APIC
  \draw[hw return] (9.4,-2.0) -- (9.4,2.0);
  \node[hw label, font=\tiny, anchor=east] at (9.25,0) {中断（MSI/引脚）};
  % DMA 数据: 设备 -> PCH -> MC
  \draw[hw bus] (9.4,-3.2) -- (9.4,-3.55) -- (1.6,-3.55) -- (1.6,-3.2);
  \node[hw label, font=\tiny] at (5.5,-3.85) {DMA 数据（直达内存，不占用 CPU 流水线）};
  \draw[hw bus] (1.6,-2.0) -- (1.6,2.0);
  \node[hw label, font=\tiny, anchor=west] at (1.75,0) {DMA 数据};
\end{pdfdiagram}
\diagramnote{控制面总图。CPU 核心居中，向外辐射三类线：实线是地址与数据路径——CPU 对 DRAM 的内存读写、对设备寄存器的 MMIO 读写、设备到内存的 DMA 数据；虚线是反方向的事件通知——设备的中断请求经中断控制器汇总后投递给 CPU。DMA 数据经根复合体直达内存控制器，途中不进入任何一条 CPU 指令的执行路径。}

\topic{写寄存器是命令，读寄存器是状态}

MMIO 写与普通内存写有本质区别。往 DRAM 地址写一个值，语义是“保存这个值，以后读回来”；往设备控制寄存器写一个值，语义是“按这个值行动”——写 \code{0x01} 可能表示启动一次发送，写 \code{0x02} 可能表示复位状态机。写的值不一定要被保存，写的动作本身才是信息。存储与互连相关章节已经指出，MMIO 写的总线响应只说明命令到达设备接口，不代表设备动作完成；从控制面看，这句话可以换个说法：MMIO 写是一条单向命令通道，完成与否要用另一条通道（状态读或中断）来确认。

MMIO 读则是对设备状态的一次采样。状态寄存器由设备硬件持续更新，CPU 读到的只是采样时刻的快照。两次读之间状态可能任意变化，这与普通内存“读到的就是上次写的”完全不同。第一节会展开这组差异如何塑造寄存器接口的各种读写约定。

\topic{中断是反方向的回信}

命令与状态两个通道都是 CPU 主动的：CPU 决定何时写、何时读。但设备完成一件事的时间点是设备才知道的，若只能靠 CPU 反复读状态去发现，CPU 就被拴在慢设备上。中断把方向反转：设备在状态变化时主动向 CPU 发一个事件信号，中断控制器负责汇总、屏蔽和排序，CPU 在指令边界接收。存储与互连相关章节第四节已经走完“事件→状态位→IRQ→控制器→向量→服务程序”的完整链路；“平台控制器：从芯片组到 IOMMU”一节会把中断控制器本身作为控制面的一台核心状态机拆开。眼下只需记住总图里的方向约定：实线由 CPU 或 DMA 发起，虚线由设备发起，控制面是双向的。

\section{二、机器里的控制理论：反馈环、稳定性与 PID}

到这里，本章已经出现了好几个控制环：VRM 把输出电压锁在参考值上，PLL 把时钟相位锁在参考沿上，固件把芯片温度维持在温度预算内、把功耗维持在功耗预算内，风扇按温度调节转速，内存训练把采样点定在眼图中心。它们的物理量、执行机构、响应速度完全不同——快的以纳秒计，慢的以秒计——但抽象成框图之后，是同一个数学结构：测量实际值，与目标比较，按差值施加纠正。这一节把这个结构本身讲清楚：先建立反馈控制的通用词汇，再回答一个任何控制环都绕不开的问题——它会不会振荡，最后用四个案例把理论读回机器里的具体电路。掌握这一节后，再遇到数据手册里“环路带宽”“补偿网络”“锁定时间”这类字样，读到的将不再是孤立的参数，而是同一个结构在不同物理载体上的表现。

\begin{keyidea}
机器里凡是要“把某个物理量维持在目标值”的地方，都有一个反馈环：传感器测量实际值，控制器按误差计算纠正量，执行器把纠正量施加到被控对象上。反馈的代价是稳定性问题——纠正信号在环路里绕行一周带回的滞后，可能让“纠正”在某个频率上变成“助推”。控制设计的全部手艺，就是在响应速度与稳定裕量之间做取舍。
\end{keyidea}

\topic{反馈控制的基本词汇}

控制理论有一套通用词汇，每个词在机器里都能找到对应的实物。先用机器里的例子把它们逐个定义清楚。

\term{被控对象}{plant}是要被维持状态的物理系统，它的某个物理量是被控量。VRM 电压环里，被控对象是 buck 变换器的 LC 输出级连同挂在上面的负载，被控量是输出电压；PLL 里，被控对象是压控振荡器，被控量是输出时钟的相位；HDD 伺服里，被控对象是磁头臂的机械结构，被控量是磁头相对磁道中心的位置；热管理环里，被控对象是硅片与散热器组成的热系统，被控量是芯片结温。

\term{传感器}{sensor}把被控量转换成可以与目标比较的电信号。VRM 用两只精密电阻组成的分压网络把输出电压按比例缩小；PLL 用鉴相器把两个时钟沿的先后变成脉冲宽度；CPU 用集成在硅片上的热敏二极管或数字温度传感器报告结温；HDD 伺服从磁盘上的伺服扇区解调出位置误差信号（存储器件相关章节第六节）。

\term{执行器}{actuator}是控制器能够推动的物理入口。buck 功率级的占空比、VCO 的控制电压、音圈电机的驱动电流、风扇的 PWM 占空比，都是执行器的输入。控制器的全部影响力都通过执行器进入物理世界，执行器的速度上限也因此决定了整个环路的响应上限。

\term{参考值}{reference, setpoint}是目标本身。VRM 的参考多由带隙基准派生——带隙基准是一个对温度和工艺都高度稳定的片上电压源（约 1.2 V），从它派生的参考电压典型在 0.6--1.25 V，经分压比折算后对应目标输出电压；PLL 的参考是外部晶体提供的时钟沿；热环的参考是写入寄存器的温度阈值。

\term{误差}{error}是参考值与测量值之差，$e = r - \hat{y}$。控制器的输入只有这一个数：它不直接知道被控对象发生了什么，只知道“偏了多少、往哪边偏”。理解反馈环的一个好习惯是追问：这个环的误差信号到底是什么物理形式？VRM 里是误差放大器两输入端的电压差，PLL 里是鉴相器输出脉冲的宽度，HDD 伺服里是 PES 数值。

最后一个区分是\term{前馈}{feedforward}与\term{反馈}{feedback}。反馈测量\emph{结果}再纠正：误差出现之后才行动，因此它能抑制模型里没有预料到的扰动，代价是慢——至少要等扰动造成可测量的偏差。前馈测量\emph{扰动}提前补偿：知道负载电流变了多少，就预先多给一份占空比；知道盘片每转的偏心规律，就按查表值预先修正磁头位置（HDD 的 RRO——repeatable runout，重复性跑偏——前馈表，见存储器件相关章节第六节）。前馈不等误差出现，因此快，但它只能补偿模型里写明的扰动，模型之外的一概无效。工程上的成熟做法是让两者分工：前馈处理已知的、可预测的扰动，反馈处理其余扰动。

\begin{pdfdiagram}[x=1cm,y=1cm]
  % 主通路：参考 -> 求和点 -> 控制器 -> 执行器 -> 被控对象 -> 输出
  \node[hw label, anchor=east] at (-0.1,0) {参考值 $r$};
  \draw[hw bus] (0.0,0) -- (0.88,0);
  \draw[BookMuted, line width=0.6pt] (1.2,0) circle (0.3);
  \node[hw label, font=\tiny] at (0.84,0.16) {$+$};
  \node[hw label, font=\tiny, anchor=west] at (1.32,-0.38) {$-$};
  \node[hw control, minimum width=1.5cm] (ctrl) at (2.9,0) {控制器};
  \node[hw compute, minimum width=1.5cm] (act) at (5.1,0) {执行器};
  \node[hw compute, minimum width=1.6cm] (plant) at (7.3,0) {被控对象};
  \draw[hw bus] (1.5,0) -- (ctrl.west);
  \node[hw label, font=\tiny] at (1.85,0.22) {误差 $e$};
  \draw[hw bus] (ctrl.east) -- (act.west);
  \draw[hw bus] (act.east) -- (plant.west);
  \draw[hw bus] (plant.east) -- (9.5,0);
  \node[hw label, anchor=west] at (9.6,0) {输出 $y$};
  \fill[BookMuted] (9.0,0) circle (0.05);
  % 扰动从上方进入被控对象
  \draw[hw danger] (7.3,1.0) -- (plant.north);
  \node[hw label, font=\tiny] at (7.3,1.28) {扰动：负载突变、温度、噪声};
  % 反馈通路：输出 -> 下 -> 传感器 -> 左 -> 上进求和点
  \node[hw state, minimum width=1.5cm] (sens) at (5.1,-1.9) {传感器};
  \draw[hw return] (9.0,0) -- (9.0,-1.9) -- (sens.east);
  \draw[hw return] (sens.west) -- (1.2,-1.9) -- (1.2,-0.32);
  \node[hw label, font=\tiny, fill=white, inner sep=1pt] at (3.0,-2.06) {测量值 $\hat{y}$};
  % 机器里的实例标注（各框下方）
  \node[hw label, font=\tiny, text width=2.6cm] at (2.9,-0.85) {误差放大器 / 环路滤波 / 伺服 PID};
  \node[hw label, font=\tiny, text width=2.4cm] at (5.1,-0.85) {buck 功率级 / VCO / 音圈电机};
  \node[hw label, font=\tiny, text width=2.4cm] at (7.3,-0.85) {LC 输出与负载 / 磁头机械结构};
  \node[hw label, font=\tiny, text width=2.4cm] at (5.1,-2.75) {分压网络 / 鉴相器 / 温度传感器};
\end{pdfdiagram}
\diagramnote{通用反馈环。实线是前向通路：参考值与测量值在求和点相减得到误差，控制器按误差计算控制量，经执行器作用于被控对象；虚线是反馈通路：传感器测量输出并送回求和点。扰动从前向通路中途进入——这是反馈存在的理由：环路无法阻止扰动到来，只能在扰动造成偏差后把它纠正回去。每个框下方给出本章各环对应的实物。}

把四个案例按这套词汇对照一遍：

\begin{longtable}{L{0.13\linewidth}L{0.40\linewidth}L{0.40\linewidth}}
\toprule
词汇 & 机制 & 机器里的实例 \\
\midrule
被控对象 & 被维持状态的物理系统与被控量 & VRM 的 LC 输出级（电压）、VCO（相位）、磁头臂（位置）、硅片与散热器（温度） \\\tablerowrule
传感器 & 把被控量变成可比较的信号 & 分压电阻网络、鉴相器、热敏二极管、伺服 burst 解调 \\\tablerowrule
执行器 & 控制器推得动的物理入口 & buck 占空比、VCO 控制电压、音圈电机电流、风扇 PWM \\\tablerowrule
参考值 & 目标值 & 带隙基准电压、参考时钟沿、温度阈值寄存器 \\\tablerowrule
误差 & $e = r - \hat{y}$，控制器的唯一输入 & 误差放大器输入差、鉴相器脉宽、PES 数值 \\\tablerowrule
前馈 / 反馈 & 测扰动提前补偿 / 测结果事后纠正 & VRM 负载线补偿；HDD 的 RRO 前馈表 + PID 闭环 \\
\bottomrule
\end{longtable}

这套词汇的价值在于可迁移。任何一个陌生的“环”，先问五个问题：被控量是什么、谁测量它、参考从哪里来、误差以什么物理形式存在、执行器推的是什么。五个问题答完，这个环的结构就清楚了，剩下的只是参数。

\topic{稳定性直觉：纠正为什么会变成助推}

负反馈在直觉上是自我纠正的：偏差出现，环路施加反向作用，偏差缩小。但任何读过示波器上电源振荡波形的人都知道，负反馈环会振荡。问题出在\emph{时间}上：纠正信号不是瞬时生效的，它在环路里每经过一个储能环节就被延迟一点。

把纠正信号在环路里绕行一周的过程拆开看。每个储能元件——电感、电容、机械惯量、热容——都引入两种效应：信号幅度被衰减，相位被推迟。在最简模型里，控制理论把每个储能环节称为一个\hwkey{极点}（LC 谐振这样的一对储能元件实为一对共轭极点），一个极点最多贡献 $-90^{\circ}$ 的相位滞后。一个极点的环最多滞后 $90^{\circ}$，纠正方向只是“打了折扣”，永远不会反；两个以上极点串联时，总滞后可以向 $180^{\circ}$ 逼近。滞后达到 $180^{\circ}$ 意味着：输出偏高时，绕环一周回来的纠正信号不是把它压低，而是恰好同相地把它推得更高——负反馈在这个频率上变成了正反馈。

此时振荡是否发生，取决于\hwkey{环路增益}：扰动绕环一周，是被放大还是被缩小。若总滞后接近 $180^{\circ}$ 的那个频率上环路增益仍大于 1，任何微小的扰动绕环一周后都变得更大，再绕一周更大——振幅持续增长，直到撞上某个物理限制（电源轨、占空比上限）变成等幅振荡。若在那个频率上增益已经小于 1，扰动每绕一周都在衰减，环是稳定的。用推秋千类比最直观：在秋千荡到最高点时推一把，能量越来越大——这是滞后凑足半个周期、推变成助推；在秋千往回荡时推，越推越小——这才是纠正。相位滞后在低频时微不足道，纠正还是纠正；频率升高后滞后累积，同一套纠正逻辑就变成了助推。

工程上把这套直觉量化成两个裕量。增益降到 1 的那个频率称为\term{穿越频率}{crossover frequency}，它近似等于环路的响应速度——带宽。\term{相位裕度}{phase margin}是穿越频率处总滞后与 $180^{\circ}$ 的距离：裕度 $45^{\circ}$--$60^{\circ}$ 是教科书式的舒适区，小于 $30^{\circ}$ 的环阶跃响应有明显振铃，趋近 $0^{\circ}$ 就是趋近持续振荡。对称地，增益裕度是滞后达到 $180^{\circ}$ 处增益与 1 的距离，习惯要求 6 dB 以上。设计一个补偿网络，本质上就是安排环路增益的下降曲线：让它在相位滞后累积到危险区之前，先降得足够低。

积分环节在这里扮演双重角色。积分器——输出是输入的累积——天然贡献恒定的 $-90^{\circ}$ 滞后，且在低频段增益很高。环里放一个积分器，稳态误差会被积成零：只要误差不为零，积分就持续累积控制量，直到误差消失为止。这是它诱人的一面；代价是那固定的 $-90^{\circ}$ 把总滞后向 $180^{\circ}$ 推近了一大步。一阶惯性环节加积分环节的组合，是机器里最常见的环路骨架：波特图上，积分器先给出每十倍频 20 dB 的下降斜率，过了惯性环节的极点频率后两段叠加，斜率加大为每十倍频 40 dB——先缓后陡；相位曲线从 $-90^{\circ}$ 起步向 $-180^{\circ}$ 滑落，补偿网络要做的就是在相位曲线滑到危险区之前，把增益曲线按到 1 以下。本节不展开波特图的作法，记住趋势就够了：\hwkey{积分换精度，带宽换速度，裕度换平稳，三者互相抢}。

\begin{classicnote}
负反馈放大器是贝尔实验室的 Harold Black 在 1927 年为解决长途电话中继放大器的失真问题发明的；Harry Nyquist 随后在 1932 年给出了稳定性的严格判据。这套理论最初的“机器”就是模拟电子线路，后来被搬到机械、化工、生物系统上，1948 年 Wiener 的《控制论》把它确立为横跨工程的通用学科。计算机里的电压环、锁相环、伺服环，直接继承的是这条电话工程的血统——VRM 误差放大器的补偿网络画法，与九十年前的中继放大器是同一张图。
\end{classicnote}

稳定性取舍在时域里最直观的表现，是阶跃响应的三种形态。给环一个突然的目标变化（或一次负载突变），输出的回应方式暴露环路的全部性格：

\begin{pdfdiagram}[x=1cm,y=1cm]
  % 坐标轴
  \draw[BookRule, line width=0.5pt, ->] (0.5,0.5) -- (10.3,0.5);
  \node[hw label, anchor=north] at (10.3,0.38) {时间};
  \draw[BookRule, line width=0.5pt, ->] (0.5,0.5) -- (0.5,4.7);
  \node[hw label, anchor=south east] at (0.5,4.7) {输出};
  % 目标值虚线
  \draw[BookMuted, dashed, line width=0.5pt] (0.5,3.4) -- (10.1,3.4);
  \node[hw label, font=\tiny, anchor=east] at (10.05,3.62) {参考值 $r$};
  % 欠阻尼：超调后振荡收敛（BookRed）
  \draw[BookRed, line width=0.85pt] (0.5,0.5) .. controls (1.2,0.6) and (1.7,3.9) .. (2.3,4.1)
    .. controls (2.9,4.3) and (3.3,2.8) .. (4.0,2.78)
    .. controls (4.7,2.76) and (5.0,3.58) .. (5.7,3.6)
    .. controls (6.4,3.62) and (6.7,3.24) .. (7.4,3.26)
    .. controls (8.1,3.28) and (8.8,3.4) .. (10.1,3.4);
  \node[BookRed, font=\tiny, anchor=west] at (4.4,4.28) {欠阻尼：超调与振铃};
  % 临界阻尼：最快无超调（BookAccent）
  \draw[BookAccent, line width=0.85pt] (0.5,0.5) .. controls (1.6,0.6) and (2.5,3.1) .. (3.7,3.34)
    .. controls (5.0,3.42) .. (7.2,3.4) -- (10.1,3.4);
  \node[BookAccent, font=\tiny, anchor=west] at (7.9,2.6) {临界阻尼：最快且无超调};
  % 过阻尼：缓慢逼近（BookTeal）
  \draw[BookTeal, line width=0.85pt] (0.5,0.5) .. controls (2.5,0.55) and (4.0,1.5) .. (5.5,2.4)
    .. controls (7.0,3.0) and (8.5,3.3) .. (10.1,3.38);
  \node[BookTeal, font=\tiny, anchor=west] at (5.6,1.85) {过阻尼：平稳但迟缓};
\end{pdfdiagram}
\diagramnote{同一阶跃目标下三种阻尼形态的响应。欠阻尼环（相位裕度小）上升最快，但冲过目标后振铃数次才停稳；过阻尼环（裕度大、带宽低）不越雷池一步，却慢得多；临界阻尼是不超调前提下最快的形态。工程上常故意留一点超调（相位裕度 $45^{\circ}$ 上下，超调约 15\%--25\%），用少量振铃换取明显更快的响应。}

三种形态对应三种设计失败或妥协。欠阻尼：补偿网络给的相位裕度太小，环在追赶目标时刹不住车。过阻尼：积分增益或带宽取得太保守，环对任何变化都慢半拍——负载突变时下陷又深又久。临界阻尼是理论最优点，实际设计宁可靠近它而偏一点欠阻尼，因为元件公差和温度会把参数推来推去，贴着临界值设计的环在批量产品中总有一部分滑进深欠阻尼区。

\begin{warningbox}
“接成负反馈就稳定”是最常见的误解。反馈的极性只在低频保证纠正方向正确；频率升高后各储能环节的相位滞后累积，负反馈会在某个频率上变成正反馈。稳定与否取决于环路增益和相位在整个频率轴上的配合，不取决于接线端的正负号。第二个常见误解是“积分增益越大、稳态越准”：积分确实把稳态误差积成零，但它的 $-90^{\circ}$ 滞后同时压缩相位裕度，积分调过头，环会以低频大幅振荡的方式“记住”这个错误。
\end{warningbox}

\topic{案例一：VRM 电压环——二阶对象与比例-积分补偿}

“主板：供电、时钟与固件”一节讲过 VRM 的结构：多相 buck 把 12 V 降到核心所需的 1 V 上下，误差放大器比较输出分压与参考，调整占空比。现在用控制理论重读这个环。

被控对象是二阶的。buck 输出级的电感与输出电容各是一个储能元件，LC 组合存在一个谐振频率（典型几十 kHz）：低于谐振频率，输出电压忠实地跟随占空比；接近谐振频率，相位急剧滑落、增益出现尖峰——若不加处理，环在谐振频率附近凑齐振荡条件。误差放大器外围的 RC 网络就是为这件事存在的：它实现\hwkey{比例-积分补偿}，积分支路在低频提供高增益，把直流输出误差积成零（这就是稳压精度的来源）；比例支路与精心放置的零点在穿越频率附近托住相位，把相位裕度维持在 $45^{\circ}$--$60^{\circ}$。补偿网络的几个电阻电容取值，就是这个环的相位裕度的物理形态。

比例-积分之外还可以加上微分环节：微分正比于误差的变化率，误差还在加速恶化时它就提前加大纠正量，相当于给环路装上预判；在频域上它提供超前相位，能部分抵消储能环节带来的滞后、抑制振荡——这正是 4.2 节“微分环节抑制振荡”的含义。只是微分对高频噪声同样敏感，会把输出纹波放大进控制量，实际 VRM 多以 PI 为主、微分为辅。

负载瞬态是观察这个环的最佳实验。CPU 从空闲跳进全负载，电流边沿的 di/dt 可达几十 A/ns，一次完整的负载阶跃在几十纳秒到微秒量级内完成（“电源、时钟与信号完整性”一章第一节算过这笔账）。环路带宽折算成响应时间是 $\mu$s 级，最初的下跌环路根本来不及参与：电压先按输出电容的等效串联电阻跌一个台阶，再因电容放电继续下滑，直到环路把占空比提上来才止步回升。下陷的深度近似为 $\Delta V \approx \Delta I \times Z_{out}$——$Z_{out}$ 是输出端看进去的等效阻抗，由电容参数与环路增益共同决定；恢复的快慢由环路带宽决定，穿越频率越高，环越早察觉、越早纠正。“电源、时钟与信号完整性”一章第一节“去耦电容按时间尺度分层”说的正是同一个事实的另一面：ns 级靠片上电容、$\mu$s 级靠板级电容，再慢才轮到环路——每一层只负责自己来得及管的时间尺度。

\begin{pdfdiagram}[x=1cm,y=1cm]
  % 负载电流阶跃
  \draw[BookAccent, line width=0.85pt] (0.5,5.0) -- (4.0,5.0) -- (4.0,6.2) -- (10.3,6.2);
  \node[hw label, anchor=east] at (0.4,5.6) {负载电流};
  \draw[BookAmber, line width=0.5pt, <->] (3.7,5.0) -- (3.7,6.2);
  \node[hw label, font=\tiny, anchor=east] at (3.55,5.6) {$\Delta I$};
  % 输出电压：基线 + 下陷 + 恢复
  \draw[BookMuted, dashed, line width=0.5pt] (4.0,3.0) -- (10.3,3.0);
  \draw[BookTeal, line width=0.85pt] (0.5,3.0) -- (4.0,3.0) -- (4.15,2.55)
    .. controls (4.4,1.95) and (4.8,1.8) .. (5.3,1.95)
    .. controls (5.8,2.1) and (6.2,3.05) .. (6.7,3.08)
    .. controls (7.2,3.11) and (7.6,3.0) .. (8.3,3.0) -- (10.3,3.0);
  \node[hw label, anchor=east] at (0.4,3.0) {输出电压};
  % 下陷深度标注
  \draw[BookAmber, line width=0.5pt, <->] (5.45,3.0) -- (5.45,1.95);
  \node[hw label, font=\tiny, anchor=west] at (6.9,2.45) {下陷 $\approx \Delta I \times Z_{out}$};
  % 恢复时间标注
  \draw[BookMuted, line width=0.5pt, <->] (4.0,1.3) -- (8.3,1.3);
  \node[hw label, font=\tiny, anchor=north] at (6.15,1.18) {恢复时间：由环路带宽决定};
  % 时间轴
  \draw[BookRule, line width=0.5pt, ->] (0.5,0.7) -- (10.3,0.7);
  \node[hw label, anchor=north] at (10.3,0.58) {时间};
\end{pdfdiagram}
\diagramnote{负载电流阶跃下的 VRM 输出电压。最初的陡跌由输出电容的寄生电阻决定，环路来不及参与；随后的下滑与回升才是环路带宽的表演时间。带宽越高的环下陷越浅、恢复越快，但带宽受 LC 谐振与相位裕度限制，不能无限推高。}

多相并联是执行器层面的优化。第四节讲过 $N$ 相 buck 错相 $360^{\circ}/N$ 驱动：各相纹波电流峰谷互错，在输出电容上叠加时大部分抵消，等效纹波频率变成 $N$ 倍、幅度大减。用控制理论的话说，这是把一个大执行器拆成 $N$ 个错峰工作的小执行器：各相电感器相同，但 $N$ 相并联后等效电感为 $L/N$，相电流对占空比的响应更快，环路可以安全地取更高的带宽；输出电容的负担也随之减轻。软件碰到这个环的方式是 4.2 节讲过的 VID 通路：CPU 把目标电压编码发给 VRM 控制器，控制器平移自己的参考值，环路自动把输出调过去。DVFS 里的“调压”没有一根模拟控制线，软件改的就是这个参考值——完整的软硬协同故事见 5.5。

\topic{案例二：PLL 与 DLL——锁相环为什么是 II 型系统}

“电源、时钟与信号完整性”一章第二节画过 PLL 的五环结构，存储器件相关章节第三节讲过 PHY 里的 DLL。用本节词汇重读：鉴相器是传感器，输出正比于相位差的脉冲；电荷泵加环路滤波器是控制器，把误差脉冲积成平缓的控制电压；VCO 是执行器，控制电压改变振荡频率；被控量是输出时钟的相位。

这个环有一个在控制理论里很特殊的性质：它天然带两个积分器。第一个是 VCO 本身——控制电压决定的是\emph{频率}，而频率是相位的导数，相位是频率对时间的积分。被控量是相位、执行器输出的是频率，这中间隔着一次免费的积分。第二个积分器在环路滤波器里：理想电荷泵把相位误差脉冲积分成控制电压。环路里串着两个积分器的系统称为 \hwkey{II 型系统}，它有两个直接后果。其一，锁定后静态相位误差为零：只要参考沿与反馈沿之间还有恒定的相位差，第二个积分器就持续累积控制电压，直到相位差消失。其二，它能无静差地跟踪频率阶跃：参考频率跳变后，环路重新锁定且相位仍然对齐。这就是为什么 PLL 是时钟倍频的标准方案——它对“准确”的承诺是结构保证的，不靠手工调参。

DLL 少一个积分器。它的执行器是数控延迟线：控制量直接就是相位（延迟），不需要经过“频率积分成相位”那一道。环路里只剩一个积分器，是 I 型系统。少一个积分器意味着少 $-90^{\circ}$ 的固有滞后，相位裕度天然宽裕，因此 DLL 锁定快、不易振铃——这正是存储器件相关章节第三节说“DLL 锁定快、引入的抖动小”的控制论原因。代价同样来自结构：延迟线不产生新频率，DLL 只能搬移相位，不能倍频。

抖动在控制理论里是环路的\hwkey{噪声传递}问题。参考时钟上的噪声进入环路后按低通特性传到输出：带宽以内的慢起伏，环路跟得上、照单全收；带宽以外的快抖动，环路来不及反应。VCO 自身的噪声正好相反，按高通特性出现：环路能纠正慢漂移，管不了快起伏。环路带宽因此是两类噪声的分界旋钮——带宽越低，参考噪声滤得越干净，但 VCO 噪声露出来的频段越宽。“电源、时钟与信号完整性”一章第二节讲“用 PLL 给抖动的参考洗干净”时把带宽设低，选的就是这个权衡的一端。软件与这个环的唯一接触点是锁定状态：启动代码先配好分频比，再轮询 locked 标志，确认环路收敛之后才切换时钟源——一个 II 型系统的收敛时间，就这样站在每台机器开机流程的关键路径上。

\topic{案例三：DVFS 与热管理——软硬协同的慢环}

前两个案例的环路带宽以 kHz--MHz 计，这一节看频谱的另一端：以秒和分钟计的热环。被控对象是硅片、导热界面材料与散热器组成的热系统：硅片功耗经热阻流向散热器，各环节的热容把温度变化拖慢——芯片结温的时间常数以秒计，散热器整体以分钟计。被控对象慢，环也快不了，而且不必快：一个热时间常数 10 秒的系统，用 1 ms 的环路去管纯属浪费。这是控制设计里一条朴素的原则：\hwkey{环路的响应速度只需匹配被控对象与扰动的速度}。

环的组成：传感器是分布在芯片热点附近的数字温度传感器（DTS）和电流监测电路；执行器有三个层次——频率（PLL 重锁或时钟分频）、电压（经 SVID 改 VRM 参考值，即 5.3 的 VID 通路）、以及功耗门控（整核时钟门控）；控制器的位置则经历了一次意味深长的迁移。早期方案（Intel SpeedStep 时代）由操作系统当控制器：OS 每隔 10--100 ms 评估一次负载，从 ACPI 定义的 P-state 表里挑一档，写 MSR 请求切换。这个环的控制周期是几十毫秒级，对编译、视频解码这类持续负载够用，但对“用户点开一个网页，渲染线程突然忙 20 ms”这种脉冲式负载完全跟不上——等 OS 察觉时任务已经完成，频率白提一次；或者更糟，升频姗姗来迟，用户先感到卡顿。Intel 从 Skylake 起把控制器搬进硬件（Speed Shift，硬件管理 P-state）：硬件以约 1 ms 的周期自主评估负载并调频，OS 退居二线，只写下一个性能偏好提示（EPP），从“每步都指挥”变成“只定基调”。控制周期从几十毫秒缩到一毫秒，环能跟踪的负载变化速度提高了一个数量级以上——这就是控制周期决定跟踪能力的直接演示。

风扇控制是同一个环在机箱尺度的延伸，结构更简单：一条\hwkey{风扇曲线}，本质是分段的\hwkey{比例控制}——温度低于阈值给固定低速，中间区间转速随温度线性上升（斜率就是比例增益），超过上限全速。纯比例控制会有稳态误差，风扇环不在乎：目标不是把温度严格保持在某一设定值，而是“避免过热并控制噪声”，误差本身被折算进了曲线设计。曲线上还有一个控制理论的经典元件：\term{迟滞}{hysteresis}。升温和降温走不同的阈值——比如升到 60$^{\circ}$C 才提速，降到 55$^{\circ}$C 才降速——否则温度在阈值附近抖动时，风扇会每秒切换一次转速，噪声反而更恼人。迟滞是用状态记忆换取决策稳定性，它是所有“阈值切换”类控制的必备零件。

\begin{lstlisting}
# 硬件监控芯片（hwmon）暴露给软件的接口：温度与风扇转速都是只读文件
$ cat /sys/class/hwmon/hwmon2/temp1_input     # 结温，单位 0.001 摄氏度
62000
$ cat /sys/class/hwmon/hwmon2/fan1_input      # 风扇实测转速 RPM
1450

; EC 固件里的风扇环（约每 100 ms 执行一次）
fan_loop:
  T = read_dts()
  if T < 50 then
      pwm = 30%                       ; 低速段（固定）
  else if T < 80 then
      pwm = 30% + (T - 50) * 2.3%     ; 比例段，斜率=增益
  else
      pwm = 100%                      ; 全速段（饱和）
  pwm = hysteresis(pwm)               ; 升降温走不同阈值
  write_pwm(pwm)
\end{lstlisting}

\begin{pdfdiagram}[x=1cm,y=1cm]
  % 坐标轴
  \draw[BookRule, line width=0.5pt, ->] (0.8,0.8) -- (10.4,0.8);
  \node[hw label, anchor=north] at (10.4,0.68) {温度};
  \draw[BookRule, line width=0.5pt, ->] (0.8,0.8) -- (0.8,4.4);
  \node[hw label, anchor=south east] at (0.8,4.4) {风扇 PWM};
  % 参考刻度
  \draw[BookMuted, dashed, line width=0.4pt] (4.0,0.8) -- (4.0,1.6);
  \node[hw label, font=\tiny, anchor=north] at (4.0,0.7) {50$^{\circ}$C};
  \draw[BookMuted, dashed, line width=0.4pt] (7.5,0.8) -- (7.5,3.9);
  \node[hw label, font=\tiny, anchor=north] at (7.5,0.7) {80$^{\circ}$C};
  \draw[BookMuted, dashed, line width=0.4pt] (0.8,3.9) -- (7.5,3.9);
  \node[hw label, font=\tiny, anchor=east] at (0.7,3.9) {100\%};
  % 风扇曲线：低速恒定 -> 比例段 -> 全速（实线为升温路径，虚线为降温路径）
  \draw[BookAccent, line width=0.9pt] (0.8,1.6) -- (4.0,1.6) -- (7.5,3.9) -- (10.2,3.9);
  \draw[BookAccent, line width=0.7pt, dashed] (0.8,1.6) -- (3.5,1.6) -- (7.0,3.9) -- (10.2,3.9);
  \draw[BookAccent, line width=0.5pt, <->] (6.0,2.91) -- (6.0,3.24);
  \node[BookAccent, font=\tiny, anchor=west] at (6.15,3.08) {迟滞};
  \node[hw label, font=\tiny, text width=3.4cm] at (9.0,2.35) {实线：升温路径；虚线：降温路径，两路径的间距即迟滞区间};
  \node[BookAccent, font=\tiny] at (2.3,1.85) {低速恒定};
  \node[BookAccent, font=\tiny, anchor=west] at (5.1,1.92) {比例段：斜率即增益};
  \node[BookAccent, font=\tiny, anchor=south] at (8.8,3.98) {全速（饱和）};
\end{pdfdiagram}
\diagramnote{风扇曲线：分段比例控制。低温段固定转速避免频繁启停；中段转速随温度线性上升，斜率就是比例增益，斜率越大控温越紧、噪声越敏感；高温段进入饱和，全速运行。50$^{\circ}$C 与 80$^{\circ}$C 两个转折点附近各有迟滞区间，防止温度抖动引起转速来回切换。}

软件在三条路径上碰到这个慢环。其一，读接口：hwmon 子系统把温度、风扇转速导出为 sysfs 文件，\code{sensors} 命令读的就是它们。其二，写接口：\code{cpufreq} 子系统的 governor（performance/powersave/schedutil）决定 OS 侧的策略，Speed Shift 平台上策略退化为一个 EPP 提示值。其三，越界接口：温度越过临界点时，硬件不等任何软件——PROCHOT 信号直接拉低频率，再不行就 THERMTRIP 强制断电。安全相关的最后一级永远做在硬件里，因为热失控时不能假设软件还活着。

\topic{案例四：内存训练——没有微分方程时的控制}

前三个案例的环有一个共同前提：被控对象可以写成连续的、近似可微的模型——电压对占空比、相位对控制电压、温度对功耗，都是光滑变化的量，误差信号连续，纠正量也连续。内存训练面对的世界完全不同。存储器件相关章节第三节讲过训练的细节：PHY 扫描每根数据线的采样延迟与判决阈值，DRAM 回报的只有 PASS/FAIL 一个比特。这里没有“误差有多大”，只有“过与没过”；眼图的形状由走线、过孔、封装、温度共同决定，无法写成一个可以求导的传递函数。连续控制的所有工具——增益、相位、补偿——在这里都无从下手。

工程上的替代方案是把“调节”换成\hwkey{系统辨识}加\hwkey{参数搜索}。第一步，测量对象：控制器把延迟和 VREF 扫成二维网格，每个格点上做一次读写测试，PASS/FAIL 填满整张表——这张表就是被控对象的特性画像，它在时间轴上的投影就是眼宽，在电压轴上的投影就是眼高。第二步，在画像上做优化：找 PASS 区域的中心点，把采样参数定在那里——中心是各方向裕量最大的位置，对日后的温度漂移和老化最稳健。搜索策略本身可以讲究：粗扫定位大区、细扫逼近边界、二分查找精确边沿，都是标准优化技术在这个离散问题上的应用。这与连续环的区别是根本性的：连续环\emph{持续地}用误差修正输出，训练是\emph{一次性地}把参数标定好，然后开环运行——它更像秤的校准，而不是定速巡航。

但这条边界正在移动。标定结果会随温度漂移：硅片的延迟随温度变化，标定时找到的眼中心，在机器跑热之后可能偏离几十 ps。LPDDR 平台与部分 DDR5 平台的应对是把一次性校准升级为慢速的持续调节：运行中周期性重训练；JEDEC DDR5 规定的则是周期性 ZQ 校准，按温度变化重新标定驱动阻抗（存储器件相关章节第三节）。训练的范式由此从“出厂一次”走向“运行中微调”，离散搜索与连续反馈在这个时间尺度上汇合了。同一趋势也出现在高速串行链路里：PCIe/USB 接收端的连续时间均衡器是模拟反馈环，判决反馈均衡器（DFE）则每个比特做一次离散调整——一条链路上两种范式并存，各管一段频率。

\topic{两种控制范式总结}

把本章和此前各章的“控制”清点一遍，可以分成两大族。一族是\hwkey{连续反馈环}：被控量是连续的物理量，误差信号连续存在，纠正量连续可调，分析工具是环路增益、相位与阶跃响应——VRM 电压环、PLL、HDD 磁头伺服（其 PID 细节见存储器件相关章节第六节，此处不再重复）、风扇与热环都在此列。另一族是\hwkey{离散状态机}：被控对象是协议与资源的状态，输入是事件与应答，纠正手段是状态转移，分析工具是状态图与时序图——内存控制器的 bank 状态机（存储器件相关章节第四节）、DMA 通道的传输状态机（存储与互连相关章节）、PCIe 枚举与链路训练、内存训练，都是这一族。

\begin{longtable}{L{0.16\linewidth}L{0.38\linewidth}L{0.38\linewidth}}
\toprule
维度 & 连续反馈环 & 离散状态机 \\
\midrule
被控对象 & 连续物理量：电压、相位、位置、温度 & 协议与资源状态：bank 行开闭、传输进度、链路状态 \\\tablerowrule
反馈信号 & 连续误差：电压差、相位差、位置误差 & 离散事件与应答：PASS/FAIL、完成中断、状态位 \\\tablerowrule
调节手段 & 连续纠正量：占空比、控制电压、电流 & 状态转移：发命令、改配置、切换模式 \\\tablerowrule
时间尺度 & 由对象时间常数决定，ns（VRM/PLL）到分钟（热） & 由协议节拍决定，一个总线周期到一次训练 \\\tablerowrule
典型失效 & 失稳：振荡、超调、锁定丢失 & 逻辑错误：死锁、超时、挂死在某个状态 \\\tablerowrule
分析工具 & 环路增益、相位裕度、阶跃响应 & 状态图、时序图、超时与重试机制 \\\tablerowrule
本章实例 & VRM 电压环、PLL/DLL、DVFS 热环、风扇 & 内存训练、bank 状态机、电源时序、链路枚举 \\
\bottomrule
\end{longtable}

两族的分界不在硬件与软件之间，而在\emph{被控对象的数学性质}之间：能写出连续模型的用环，写不出的用状态机。真实的控制器几乎总是两者的复合体——VRM 芯片里，电压环是连续反馈，上电时序、过流保护、多相切相是状态机；内存控制器里，调度是状态机，DLL 与周期性重训练是环；HDD 里，寻道是状态机切换三段轮廓，磁道跟踪是连续伺服。整机之所以可靠，不是因为哪一族更优越，而是每个问题都被交给了与它数学性质匹配的那一族。

读懂一台机器里所有的控制环之后，“CPU 如何控制整台机器”这个问题就有了第二层答案：CPU 并不直接控制机器，它配置的是一大批自治的控制环与状态机——设定参考值、读取状态字、接收中断——真正的控制在环里，以各自的物理节奏持续进行，无论 CPU 此刻在不在看它们。
\section{三、软件侧的控制：驱动、中断与调度}

本章到目前为止讨论的控制环都运行在硅片和固件里：VRM 在纳秒尺度上稳住电压，PLL 在皮秒尺度上对齐相位，内存控制器在每个时钟周期决定发哪条命令。这些环对软件完全透明——操作系统看不到它们，也不该看到。但整台机器还有另一半控制问题：设备谁来找、找到了谁来管、设备有事了谁被叫醒、CPU 时间按什么规则分。这一层的控制周期在微秒到秒之间，控制器就是操作系统本身。本节回答三个问题：操作系统怎么发现设备并把驱动的代码接到设备寄存器上，设备事件怎么沿中断链一路上行到某个等待的进程，以及定时器中断怎样把调度器变成一个离散时间控制器。

\begin{keyidea}
软件侧的控制和硅片里的控制环是同一套结构，只是时间尺度不同。设备寄存器是执行器的接口，中断是传感器送到 CPU 的信号，驱动是环路的控制律实现。把它们当作控制环来读，驱动代码里“读状态—发命令—等完成”的固定套路就不是经验习惯，而是闭环的必要环节。
\end{keyidea}

\topic{操作系统怎么发现设备}

驱动要能控制设备，前提是操作系统先知道设备存在、知道它的寄存器在地址空间的哪个窗口。发现设备有三条途径，分别对应三种硬件现实。

第一条途径是 \term{ACPI 表}{Advanced Configuration and Power Interface}。PC 和服务器上，固件在内存里放一组表，其中 DSDT 等表用 ACPI 机器语言描述主板上“焊死”的设备——嵌入式控制器、温度传感器、电源按钮——以及每个设备占用的 I/O 端口、MMIO 窗口和中断线。操作系统启动时解析这套命名空间，为每个设备节点建立内部对象，再按节点里的硬件标识符（HID）去驱动库里找匹配的驱动。第二条途径是\term{设备树}{device tree}。嵌入式 ARM、RISC-V 平台没有固件代为描述的传统，改为 bootloader 把一份二进制设备树（DTB）随内核一起传入内存，里面逐节点列出 SoC 内部总线和挂在上面的控制器；驱动按节点里的 compatible 字符串声明“我能管这类设备”，内核按字符串配对。第三条途径是 \term{PCI 枚举}{PCI enumeration}：PCI/PCIe 设备是可插拔的，无法事先写进任何静态表，但协议本身提供了发现手段——每个设备的开机配置空间里都有厂商号、设备号，主机按总线号逐设备读出这些编号，再读出 BAR 寄存器声明的 MMIO 窗口大小，把窗口分配进物理地址空间。驱动携带一张自己支持的编号表，枚举到一个设备就比对一次，命中即接管。经典芯片与平台案例讲过平台启动时枚举在开机流程中的位置，这里关心的是三条途径的共同产物：一个“设备对象 + 一组已分配的地址窗口 + 一段被调用的驱动入口函数”。

\begin{longtable}{L{0.12\linewidth}L{0.32\linewidth}L{0.20\linewidth}L{0.27\linewidth}}
\toprule
途径 & 信息载体 & 典型平台 & 驱动匹配的入口 \\
\midrule
ACPI 表 & 固件放置在内存中的表，描述设备存在与资源分配 & PC、服务器 & 解析命名空间，按硬件标识符（HID）找驱动 \\\tablerowrule
设备树 & bootloader 随内核传入的二进制描述（DTB） & 嵌入式 ARM、RISC-V & 按 compatible 字符串配对驱动 \\\tablerowrule
PCI 枚举 & 总线协议自身：配置空间里的厂商号、设备号与 BAR & 所有含 PCI/PCIe 的平台 & 按驱动自带的编号表逐项比对 \\
\bottomrule
\end{longtable}

三条途径的分工由硬件的可插拔性决定：焊在主板上的设备写进静态描述（ACPI 或设备树），插槽上的设备靠协议自己报告（PCI 枚举）。匹配成功之后，驱动拿到的是设备窗口的物理地址，还不能直接用——x86-64 内核运行在虚拟地址上，设备窗口必须先映射进来。\code{ioremap} 完成这件事：它在内核虚拟地址空间中建立一段到该物理窗口的映射，并给这段映射打上设备内存属性（不可缓存、不可合并、写顺序受保护），这些属性正是存储与互连相关章节反复强调的 MMIO 访问约定。

映射建立之后，驱动读写设备寄存器和读写普通内存的代码在指令层面完全一样——都是 \code{mov}。下面是一个最小设备驱动的典型骨架，左侧是 C 形式，右侧是它编译后对应的 x86-64 指令形态：

\begin{lstlisting}
/* base 是 ioremap 返回的内核虚拟地址；
 * SR: 状态寄存器 (bit0 = BUSY)，CR: 命令寄存器，IE: 中断使能 */
while (readl(base + SR) & BUSY)      /* 读 SR 判忙 */
    cpu_relax();
writel(CMD_START, base + CR);        /* 写 CR 发命令 */
writel(IRQ_DONE,  base + IE);        /* 开中断，然后睡眠等待 */

/* 上面三行编译后就是普通的 mov，与读写内存同一类指令：
 * mov  eax, [rbx]          ; 读 SR
 * test eax, 1
 * jnz  .busy
 * mov  dword [rbx+4], 1    ; 写 CR
 * mov  dword [rbx+8], 1    ; 写 IE                          */
\end{lstlisting}

代码和内存访问长得一样，语义却完全不同：读状态寄存器拿到的是设备此刻的硬件事实，写命令寄存器是推了一下设备内部的状态机。存储与互连相关章节已经指出关键一点——MMIO 写返回只说明命令到达设备接口，不代表动作完成。完成的通知有两条路：驱动死等状态位（轮询），或者像上面这段代码一样开好中断去睡眠，把 CPU 让给别的进程。慢设备上轮询是纯粹的浪费，于是需要中断把“完成”这件事主动送回来——这是下一小节的主题。

\topic{中断处理链：从设备事件到被唤醒的进程}

存储与互连相关章节从硬件角度讲过中断链路：事件在设备内锁存成状态位并拉起请求线，中断控制器按屏蔽位和优先级裁决，CPU 在指令边界采样、保存现场、查表跳转。这里把同一条链按现代 x86-64 平台补全，并讲清链路的软件后半段。

现代平台上有两种中断投递方式。老方式是引脚：设备拉起一条物理中断线，进入 \term{IO/APIC}{I/O Advanced Programmable Interrupt Controller}——现代平台上它集成在 PCH 内（早期主板上是独立的 82093AA 芯片），\hwevidence{典型实现有 24 条重定向输入}，每条对应一张重定向表项，写着这条线应翻译成几号向量、投给哪个核。新方式是 \term{MSI}{message signaled interrupt}：设备不拉线，而是向一段特殊地址窗口做一次内存写事务，写的内容里直接带着向量号和目标核。MSI 省掉了引脚和共享线的竞争，还能让一个设备拥有多个向量——每核一个队列一个向量，“外设、网络与人机接口”一章里 NVMe 的每核发中断靠的就是它。两条路殊途同归，终点都是目标核的\term{本地 APIC}{Local APIC}：它按优先级决定何时把向量递交给 CPU。

CPU 在指令边界接收向量后做的事由硬件完成：保存返回地址和标志，按向量号查\term{中断描述符表}{interrupt descriptor table, IDT}——一张 256 项的表，每项是一个处理函数的入口——跳到对应表项。从这里开始是软件。处理函数即 \term{ISR}{interrupt service routine, 中断服务例程}，它最急的工作只有两件：按设备规则清掉中断源（否则同一事件会反复触发），以及把设备里等着的少量数据取到内存安全的地方。做完这两件，ISR 就应该退场，把剩下的慢活——协议解析、缓冲区整理、唤醒逻辑——推迟给下半部：\term{软中断}{softirq}、\code{tasklet} 或者内核线程，由它们在允许睡眠、允许被抢占的上下文里慢慢做完，最后唤醒睡眠在等待队列上的进程。整条链如下：

\begin{pdfdiagram}[x=1cm,y=1cm]
  % --- 主链：设备 → IO/APIC → Local APIC → CPU → ISR ---
  \node[hw compute, minimum width=1.7cm] (dev) at (0,0) {设备\\事件/状态位};
  \node[hw control, minimum width=1.7cm] (ioapic) at (2.5,0) {IO/APIC\\重定向表};
  \node[hw control, minimum width=1.7cm] (lapic) at (5.0,0) {Local APIC\\IRR/ISR};
  \node[hw compute, minimum width=1.7cm] (cpu) at (7.5,0) {CPU 核\\查 IDT};
  \node[hw control, minimum width=1.7cm] (isr) at (10.0,0) {ISR 上半部\\清源/取数};
  % --- 下半部与进程（下方一行） ---
  \node[hw compute, minimum width=1.7cm] (bh) at (10.0,-2.2) {下半部\\softirq/线程};
  \node[hw state, minimum width=1.7cm] (proc) at (7.5,-2.2) {等待进程\\被唤醒};
  % --- 主链箭头 ---
  \draw[hw bus] (dev) -- node[above,hw label]{中断线} (ioapic);
  \draw[hw bus] (ioapic) -- node[above,hw label]{向量} (lapic);
  \draw[hw bus] (lapic) -- node[above,hw label]{投递} (cpu);
  \draw[hw bus] (cpu) -- node[above,hw label]{保存现场} (isr);
  \draw[hw bus] (isr) -- node[left,hw label,fill=white,inner sep=1pt]{推迟慢活} (bh);
  \draw[hw bus] (bh) -- node[below,hw label]{唤醒} (proc);
  % --- MSI 直通路径：设备直接写消息到 Local APIC ---
  \draw[hw bus] (dev.south) -- (0,-1.1) -- (5.0,-1.1) -- (lapic.south);
  \node[hw label, fill=white, inner sep=1pt] at (2.5,-1.1) {MSI：向 APIC 地址写消息};
  % --- EOI 返回路径 ---
  \draw[hw return] (isr.north) -- (10.0,1.1) -- (5.0,1.1) -- (lapic.north);
  \node[hw label, fill=white, inner sep=1pt] at (7.6,1.1) {EOI：恢复投递};
\end{pdfdiagram}
\diagramnote{中断处理链。实线是事件上行路径：引脚方式经 IO/APIC 重定向，MSI 方式由设备直接向 Local APIC 的地址窗口写消息；CPU 按向量查 IDT 进入 ISR，上半部只做清源和取数，慢活推迟给下半部，最后由下半部唤醒等待的进程。虚线是返回路径：ISR 向 Local APIC 发 EOI，控制器才恢复同级和更低优先级请求的投递。}

\begin{warningbox}
“中断处理函数里可以慢慢处理”是最常见的驱动错误。ISR 运行在中断上下文：它没有自己的进程身份，不能睡眠——一旦睡眠，没有任何东西能把它再调度回来；它还可能正压着同级和更低优先级的中断（EOI 之前这些请求都在等待）。在 ISR 里做耗时操作，后果不是这个设备变慢，而是整台机器的中断延迟被拉长，别的设备开始丢数据。规则只有一条：上半部做最急的事，其余全部推迟。
\end{warningbox}

\topic{定时器中断驱动调度：操作系统的心跳}

中断链上最特殊的一个“设备”是定时器：它不报告外部事件，而是周期性地打断 CPU，给操作系统一个固定的决策节拍。现代 x86-64 平台上，这个节拍来自每个核本地的 \term{Local APIC 定时器}{Local APIC timer}——它按可编程的周期产生中断，内核把它配置成每秒 \hwtiming{250 或 1000 次}（内核配置项 HZ），即每 \hwtiming{4 ms 或 1 ms} 一个节拍，称为一次 \hwkey{tick}。

每个 tick 的处理路径很短，做的却是控制环里“传感器”和“决策”两件事。传感器一侧：更新系统计时，并给当前正在运行的进程累计这一段实际运行时间。决策一侧：检查当前进程的时间片是否用完、是否有更高优先级的进程变得可运行；需要换人就设置一个重调度标志，在中断返回用户态之前检查这个标志、进入调度器。调度器选中下一个进程，执行上下文切换——把当前进程的寄存器现场存进它的进程描述符，把下一个进程的现场装回 CPU。

把这个结构用控制语言写出来，它就是一台教科书式的离散时间控制器：被控量是各进程实际获得的 CPU 时间分配，传感器是定时器中断和运行时间累计，控制器是调度算法，执行器是上下文切换。与硅片里的连续环不同，它按固定节拍采样、按节拍决策——控制周期就是 tick。调度算法本身以 CFS（Completely Fair Scheduler，完全公平调度器）为例：它给每个进程维护一个虚拟运行时间 vruntime，把实际运行时间按优先级加权折算后累加，每次调度选 vruntime 最小的进程运行，使长期来看各进程的 CPU 份额收敛到公平值。至此，整机控制链的最后一环闭合了：连“CPU 这台执行器把时间分给谁”，也是一个有传感器、有控制律、有执行器的反馈环。

\begin{pdfdiagram}[x=1cm,y=1cm]
  \node[hw control, minimum width=1.75cm] (timer) at (0,0) {Local APIC\\定时器};
  \node[hw compute, minimum width=1.75cm] (tick) at (2.7,0) {tick 处理\\计时/累计};
  \node[hw control, minimum width=1.75cm] (sched) at (5.4,0) {调度器\\比较/选择};
  \node[hw compute, minimum width=1.75cm] (ctxsw) at (8.1,0) {上下文切换\\（执行器）};
  \node[hw state, minimum width=1.75cm] (proc) at (10.8,0) {当前进程\\（被控对象）};
  \draw[hw bus] (timer) -- node[above,hw label,font=\tiny]{节拍中断} (tick);
  \draw[hw bus] (tick) -- node[above,hw label,font=\tiny]{调度标志} (sched);
  \draw[hw bus] (sched) -- node[above,hw label,font=\tiny]{选下一个} (ctxsw);
  \draw[hw bus] (ctxsw) -- node[above,hw label,font=\tiny]{切换} (proc);
  \draw[hw return] (proc.south) -- (10.8,-1.2) -- (2.7,-1.2) -- (tick.south);
  \node[hw label, font=\tiny, fill=white, inner sep=2pt] at (6.75,-1.2) {实测运行时间：每个 tick 累计给当前进程};
\end{pdfdiagram}
\diagramnote{调度器作为离散时间控制器。Local APIC 定时器以固定节拍产生中断（采样）；tick 处理完成计时并把这一段运行时间累计给当前进程（反馈）；调度器比较后选中下一个进程；上下文切换把决策施加到 CPU 上。控制周期等于 tick 间隔，被控量是各进程实际获得的 CPU 时间。}

\topic{边界原则：哪些环必须在硬件里，哪些可以放软件}

到这一步，本章出现过的控制环已经分布在从皮秒到秒的十二个数量级上。它们各自落在硬件、固件还是软件里，不是历史偶然，而是一条可计算的边界。判据有两条。第一条来自控制理论：\hwkey{控制周期必须比被控对象的时间常数快至少一个量级}，否则扰动已经发展完成，纠正动作才刚到达，环路不是稳住对象而是跟着振荡。第二条来自软件的物理下界：软件环最短的周期是一次中断加一次上下文切换，典型在微秒量级；再慢的策略性软件（用户态守护进程）响应在毫秒到秒。凡是被控对象的时间常数小于这个下界的环，没有任何软件选项，只能放在硅片里。

\begin{longtable}{L{0.22\linewidth}L{0.24\linewidth}L{0.20\linewidth}L{0.26\linewidth}}
\toprule
控制环 & 被控对象的时间常数 & 控制周期（量级） & 实现位置 \\
\midrule
VRM 电压环 & 负载瞬变，纳秒级 & 微秒级 & 硬件：VRM 控制器（模拟环） \\\tablerowrule
PLL / DLL 相位环 & 相位漂移，皮秒至纳秒级 & 连续（锁定微秒级） & 硬件：鉴相器、VCO、延迟线 \\\tablerowrule
DRAM 命令调度 & 总线空闲窗口，纳秒级 & 每时钟周期 & 硬件：内存控制器 \\\tablerowrule
刷新调度 & 电容保持，毫秒级 & \hwevidence{7.8 $\mu$s}（高温 3.9 $\mu$s） & 硬件：内存控制器定时 \\\tablerowrule
HDD 伺服 & 机械运动，kHz 级带宽 & 数十 kHz 采样 & 硬件+盘控固件 \\\tablerowrule
TFC 飞行高度 & 加热器热惯性，毫秒级 & kHz 至 ms 级 & 盘控固件 \\\tablerowrule
I/O 调度 & 请求到达间隔，微秒至毫秒级 & 微秒至毫秒级 & 软件：内核块层 \\\tablerowrule
DVFS & 负载变化，毫秒至秒级 & 毫秒级 & 软件/固件：调频策略 \\\tablerowrule
热管理、风扇 & 整机热惯性，秒级 & 百毫秒至秒级 & 软件/固件：EC、OS \\
\bottomrule
\end{longtable}

\begin{classicnote}
“采样频率远高于被控对象带宽”是离散控制的通用前提，与信号重建里的奈奎斯特定理在精神上相通：你要控制一个对象，就得比它变得快。工程上取十倍以上的频率裕量，才有空间放下控制律本身的延迟。把这条判据反过来读，就得到整机设计的一条分工线——比微秒快的一切环必须进硅片，比毫秒慢的一切环都可以写成代码，中间地带（kHz 级的机械环）属于固化在设备里的固件。
\end{classicnote}

这条边界也解释了前几节的安排：存储器件相关章节第四节里 FR-FCFS 调度必须做成硬件状态机，因为 DRAM 总线的空闲窗口只有几纳秒，任何软件都赶不上；而同样是“调度”，块层的 I/O 调度器可以从容地用微秒级软件实现，因为它面对的对象（队列深度、请求到达）本身就慢三个量级。环放在哪里，是被控对象的速度决定的，不是实现者的偏好。

\section{四、整机控制地图}

本书从“状态、时钟与时序逻辑”一章的交叉耦合反相器开始，一路讲到的几乎每一个子系统，内部都至少有一个反馈环。现在把它们全部列在一张表里：环路的被控量、传感器、执行器、控制器类型、时间尺度，以及它在本书哪里讲过。这张表是本章的收尾，也是全书硬件部分的一张索引。

\tablechunkintro{1}{2}{VRM 电压环}{HDD 伺服}
\begin{longtable}{L{0.13\linewidth}L{0.11\linewidth}L{0.13\linewidth}L{0.11\linewidth}L{0.12\linewidth}L{0.09\linewidth}L{0.17\linewidth}}
\toprule
环路 & 被控量 & 传感器 & 执行器 & 控制器类型 & 时间尺度 & 讲过的地方 \\
\midrule
VRM 电压环 & 核心电压 & 输出端分压采样 & 功率级占空比 & 模拟 PI / 滞环 & ns--$\mu$s & “电源、时钟与信号完整性”一章 \\\tablerowrule
PLL 相位环 & 时钟相位与频率 & 鉴相器 & VCO 控制电压 & 模拟锁相环 & ps--ns & “电源、时钟与信号完整性”一章 \\\tablerowrule
DLL 延迟环 & 时钟延迟 & 相位比较器 & 可变延迟线 & 数字延迟锁定 & ps--ns & 存储器件相关章节第三节 \\\tablerowrule
内存训练 & 采样点位置 & 参考图案读回比对 & 延迟线、驱动强度、ODT & 固件扫参状态机 & 一次性，ms & 存储器件相关章节第三节 \\\tablerowrule
DRAM 命令调度 & 总线利用率与延迟 & 请求队列、bank 状态 & 命令发射顺序 & FR-FCFS 硬件调度 & ns & 存储器件相关章节第四节 \\\tablerowrule
刷新调度 & 电容电荷保持 & 定时器、温度传感器 & REF 命令 & 硬件定时加温度补偿 & $\mu$s & 存储器件相关章节第二节 \\\tablerowrule
HDD 伺服 & 磁头位置（PES） & 伺服扇区 burst & 音圈电机电流 & 数字伺服、双级执行器 & kHz & 存储器件相关章节第六节 \\
\bottomrule
\end{longtable}

\tablechunkintro{2}{2}{TFC 飞行高度}{风扇与热管理}
\begin{longtable}{L{0.13\linewidth}L{0.11\linewidth}L{0.13\linewidth}L{0.11\linewidth}L{0.12\linewidth}L{0.09\linewidth}L{0.17\linewidth}}
\toprule
环路 & 被控量 & 传感器 & 执行器 & 控制器类型 & 时间尺度 & 讲过的地方 \\
\midrule
TFC 飞行高度 & 磁头--盘片间距 & 读信号幅度 & 滑块加热器功率 & 盘控固件闭环 & ms & 存储器件相关章节第六节 \\\tablerowrule
cache MSHR（未命中状态保持寄存器） & 在途 miss 数 & miss 计数、队列占用 & 合并、阻塞、回退 & 硬件队列管理 & ns & 存储与互连相关章节、存储器件相关章节第九节 \\\tablerowrule
NVMe 队列流量控制 & 在途命令数 & 完成队列计数 & 提交节奏、门铃 & 指针差限深 & $\mu$s--ms & 存储与互连相关章节、“外设、网络与人机接口”一章 \\\tablerowrule
OS 调度 tick & CPU 时间分配 & tick 中断 + 计时 & 上下文切换 & CFS & ms & “软件侧的控制：驱动、中断与调度”一节 \\\tablerowrule
DVFS & 性能与功耗 & 负载计数器、功耗计 & 频率与电压档位 & OS/固件调频策略 & ms--s & “软件侧的控制：驱动、中断与调度”一节 \\\tablerowrule
风扇与热管理 & 芯片温度 & 片内热敏二极管 & 风扇 PWM、降频 & 滞环 / PID & s & 本章、经典芯片与平台案例 \\
\bottomrule
\end{longtable}

把这张表按时间尺度从上往下读，一个结构浮现出来：整台机器就是几十个控制环的嵌套，时间尺度从皮秒一直排到秒。快的环离物理近——VRM 贴着功率级，PLL 贴着时钟树，DRAM 调度贴着命令总线，它们必须做进硅片，因为被控对象的变化比任何软件的反应都快；慢的环离物理远——DVFS 看的是毫秒级的负载趋势，风扇看的是秒级的温度走向，调度器看的是毫秒级的公平性，它们可以放在固件和操作系统里，用可升级的代码实现。中间地带属于设备自己的固件：硬盘伺服和 TFC 运行在盘控处理器上，内存训练运行在开机固件里——它们比软件快、又不必像模拟环那样连续，恰好落在可编程但不可打扰的位置。

这张地图同时回答了本书开头隐含的那个问题：什么叫理解一台机器。一半答案是数据通路——数据从哪里来、经过哪些变换、到哪里去，从取指到执行，从内存层次到 DMA 回填；另一半答案是控制面——谁按什么规则让这一切发生：哪个环稳住电压，哪个环对齐相位，哪个环决定下一条命令，哪个环分配时间。数据通路决定机器能做什么，控制面决定机器在什么条件下持续地做下去。两者都看见了，机器才不再是黑箱。
\section{五、RAS：错误检测、约束与恢复}

RAS 分别关注可靠性、可用性与可维护性。可靠性强调在规定时间内给出正确结果；可用性强调出现局部故障后系统仍能提供服务；可维护性强调能否定位、更换并验证故障部件。三者并不总是一致：立即停机能保护数据正确性，却降低当下可用性；忽略可纠正错误能减少告警，却可能错过器件劣化趋势。

硬件错误首先按影响分为 corrected、deferred 与 uncorrected。可纠正错误由 ECC 或重试恢复，当前操作可继续，但应记录位置和频率；deferred 错误暂未影响当前指令，却表明某份数据已经损坏，必须在未来使用前隔离；不可纠正错误若已污染架构状态，则需要机器检查、任务终止或系统复位。报告“发现错误”不等于已经阻止错误数据被消费，必须继续追踪传播边界。

\begin{longtable}{L{0.18\linewidth}L{0.31\linewidth}L{0.41\linewidth}}
\toprule
机制 & 能力边界 & 仍需配套 \\
\midrule
奇偶校验 & 检出部分奇数位错误，通常不能纠正 & 重试、失效或错误上报 \\\tablerowrule
SECDED ECC & 纠正单 bit、检测双 bit（具体能力依编码而定） & scrub、地址记录与坏页隔离 \\\tablerowrule
CRC/LCRC & 检测传输中突发错误 & 序号、重放缓冲与超时 \\\tablerowrule
watchdog & 检测控制流长时间无进展 & 安全复位域、状态保留和原因日志 \\\tablerowrule
冗余副本 & 在副本独立失效时恢复服务 & 一致性、仲裁和共同故障域分析 \\
\bottomrule
\end{longtable}

\section{六、错误如何穿过整机}

一次内存 ECC 错误的路径可能是：内存控制器检测综合征；若可纠正则修正返回数据并递增计数；把物理地址、通道、DIMM、bank 和 syndrome 写入错误寄存器；中断或机器检查通知固件与操作系统；操作系统聚合同一页的错误频率，必要时迁移数据并下线页面；管理软件最终把部件位置映射为可更换单元。

这条路径中的寄存器必须考虑“首次错误”和“后续错误”。若后续事件覆盖首次地址，根因会丢失；若只锁存首次事件又不提供计数，错误风暴会被低估。常见做法是首错锁存、分类计数、溢出标志和明确清除协议，并在复位后保留足以解释上一次关机的记录。

互连错误同样需要闭合。请求可能在源端成功发送，却因目标不存在、权限错误、CRC 重试耗尽或目标超时而失败。响应要携带可归属的 requester/tag，使异常能够回到发起指令、DMA 队列或驱动请求；无法归属的错误则升级为平台级事件。仅在控制台打印“总线错误”不能支持恢复。

\section{七、复位域、隔离与降级运行}

整机复位不是唯一恢复手段。设计可把 CPU 核、PCIe function、控制器、PHY 或外设划成独立复位域，在保留其他服务的同时恢复局部状态。但复位边界必须覆盖所有在途事务和缓存状态：停止新请求，等待或取消旧请求，屏蔽中断，保存诊断信息，执行复位，重新训练与枚举，重建队列，最后恢复流量。

若设备支持功能级复位，驱动仍不能假定内存中的旧描述符自动失效。复位前后 generation 或 queue phase 必须变化，使迟到完成不会被当作新请求；IOMMU 映射和 DMA 所有权也要在设备确实停止后才能撤销。

降级运行要明确性能和安全边界。例如内存通道下线会改变交织映射，链路降速会改变带宽与超时假设，风扇故障后的功耗限制会降低处理能力。控制面应把这些状态提供给调度与运维，而不是继续使用正常容量和时延模型。

\section{八、遥测、故障注入与可验证恢复}

可维护系统应让每个重要队列观察提交、完成、超时与重置次数；让每个链路观察纠错、重放、降速和失步；让每个存储层观察校验失败、介质重试与持久化屏障。时间戳要说明来自哪个时钟域，计数器要说明清零和饱和行为。跨层关联通常依赖请求标识、物理地址、设备 BDF/队列号和软件 trace id。

恢复流程若从不被触发，就不能证明可用。故障注入可以在受控环境中制造 ECC、丢中断、完成超时、链路复训和突然断电，验证系统是否在规定边界内检测、隔离、恢复并保留证据。注入点要位于真实检测路径之前，且测试结束后确认所有强制错误开关已经清除。

\begin{pdfdiagram}[
  node distance=8mm and 8mm,
  rasbox/.style={draw, rounded corners, align=center, minimum width=24mm, minimum height=9mm, fill=blue!5},
  >=Stealth]
\node[rasbox] (detect) {检测与纠正};
\node[rasbox, right=of detect] (record) {锁存证据\\分类计数};
\node[rasbox, right=of record] (contain) {阻止传播\\隔离故障域};
\node[rasbox, below=of contain] (recover) {重试、复位\\或降级};
\node[rasbox, left=of recover] (verify) {验证服务与\\数据一致性};
\node[rasbox, left=of verify] (trend) {趋势分析与\\维护决策};
\draw[->] (detect) -- (record);
\draw[->] (record) -- (contain);
\draw[->] (contain) -- (recover);
\draw[->] (recover) -- (verify);
\draw[->] (verify) -- (trend);
\draw[->] (trend) -- node[left, font=\small]{改进阈值与策略} (detect);
\end{pdfdiagram}
\diagramnote{平台错误从检测到恢复的闭环。恢复成功仍要验证数据与服务状态，并把错误趋势用于后续阈值和维护决策。}




























\section{九、从 ECC 综合征走到可执行的错误策略}

“有 ECC”只说明编码器与检查器存在，平台仍要定义综合征如何分类、数据是否可交付、证据写到哪里以及何时升级恢复。以扩展 Hamming SECDED 为例，普通 Hamming 校验给出 syndrome，额外总奇偶位给出 overall parity。四种组合决定不同动作。

\begin{longtable}{L{0.18\linewidth}L{0.18\linewidth}L{0.26\linewidth}L{0.28\linewidth}}
\toprule
syndrome & overall parity & 解释 & 数据动作 \\
\midrule
0 & 0 & 未检测到错误 & 正常交付 \\\tablerowrule
非 0 & 1 & syndrome 指向的单 bit 错误 & 翻转该 bit 后交付并记录 \\\tablerowrule
0 & 1 & 额外总奇偶位自身出错 & 修正校验位并记录 \\\tablerowrule
非 0 & 0 & 检测到双 bit 错误 & 不可按 syndrome 误纠正 \\
\bottomrule
\end{longtable}

例如 syndrome 为 6 且 overall parity 为 1，可把码字第 6 位翻转后交付；若 syndrome 仍为 6 而 overall parity 为 0，则符合双 bit 错误模式，不能把第 6 位改掉后声称数据已经恢复。实际内存控制器还会交织多个器件与通道，syndrome 到“哪根 DQ、哪颗 DRAM”的映射必须结合布线和控制器配置解释。

\topic{scrub 周期是容量、带宽与风险的共同预算}

patrol scrub 周期性读取内存，使潜伏单 bit 错误在遇到第二个错误前被纠正并写回。若可用内存为 256 GiB，scrub 带宽为 1 GiB/s，一次完整扫描至少需 256 s；提高到 8 GiB/s 可降到 32 s。若内存可用带宽按 100 GiB/s 估算，1 GiB/s scrub 约占 1\%，但实际干扰还取决于 bank 冲突、刷新和业务局部性。

周期太长会增大两个独立错误在同一码字累积的窗口，周期太短又会抢占带宽、增加功耗。平台应记录实际完成周期和被业务背压的时间，而不是只保存配置目标。发现已纠正错误后可以触发 demand scrub 或页面迁移，但要限制错误风暴下的额外流量，防止恢复动作本身压垮系统。

\section{十、遏制边界决定错误能否从局部升级为全局}

错误检测位置不一定就是污染边界。Cache 读出不可纠正数据时，可以把对应 line 标为 poison，阻止它被当成普通数据静默使用；互连把 poison 与请求身份带到消费端，操作系统再根据地址归属决定杀死进程、下线页面或升级机器检查。若 poison 位在协议桥处丢失，硬件虽然“检测过”，错误数据仍可能被设备或 CPU 使用。

\begin{longtable}{L{0.16\linewidth}L{0.24\linewidth}L{0.28\linewidth}L{0.22\linewidth}}
\toprule
边界 & 必须保留的身份 & 可选遏制动作 & 升级条件 \\
\midrule
Cache/内存 & 物理地址、syndrome、读写类型 & 纠正、poison、失效 line & 数据已进入不可撤销状态 \\\tablerowrule
互连 & requester、事务 ID、目标 & 错误响应、隔离端口 & 无法归属或协议状态失真 \\\tablerowrule
I/O 设备 & BDF、PASID、队列、generation & 停队列、FLR、撤映射 & DMA 仍无法停止 \\\tablerowrule
操作系统 & 页面、进程、文件或虚机所有者 & 下线页、杀任务、迁移服务 & 共享内核状态受污染 \\
\bottomrule
\end{longtable}

deferred error 的意义正是“现在尚未消费，但未来使用前必须处理”。平台若只发一次日志却仍允许该页面被重新分配，就把确定的潜伏错误变成不确定的后续崩溃。错误状态要随对象生命周期移动：页面迁移后旧页下线，设备复位后旧 generation 拒绝，虚机迁移后错误归属仍能追到原工作负载。

\topic{阈值策略要有时间窗、迟滞和有限状态}

单次 corrected error 可以记录后继续，重复错误则可能预示器件劣化。阈值不能只写“超过 N 次告警”，还要说明对象和时间窗。例如同一物理页 24 小时内出现 4 次 corrected error 时迁移并下线；同一 DIMM 一小时内有 100 个不同页面报告错误时升级部件维护。前者捕捉局部弱单元，后者捕捉范围性问题。

状态机可分为 healthy、watch、degraded 和 isolated。进入高风险状态用较高阈值，恢复到低风险状态用更低阈值并要求保持一段安静期，形成迟滞，避免温度或负载在边界附近抖动时反复下线和上线。所有阈值都要能回放：输入计数、窗口起点、状态转移原因和执行动作必须进入日志。

\section{十一、共同故障域与恢复验收}

两个副本只有在失效足够独立时才提供预期冗余。双电源若接在同一 PDU，双网卡若共用同一交换机，两个固件槽若由同一个易错更新器同时改写，都有共同故障域。可靠性图不能只数部件数量，还要沿电源、时钟、冷却、固件、布线和运维动作寻找共享依赖。

恢复成功也不能只看“设备重新出现”。一块网卡复位后至少要验证：旧 DMA 不再发生，IOMMU 映射与新队列一致，旧 completion 被 generation 拒绝，中断向量只对应新队列，收发序号和统计没有重复或遗漏，上层连接按策略恢复。平台级故障注入的验收应写成不变量，而不是观察到一次成功 ping 就结束。

\begin{keyidea}
完整 RAS 闭环是：检测并分类，保存首错与趋势证据，阻止错误跨过遏制边界，选择重试/复位/降级，最后验证数据与服务不变量。任何“自动恢复”若缺少最后一步，都只能证明状态机执行过，不能证明系统恢复正确。
\end{keyidea}

\section*{章末练习}

\begin{chapterexercise}{1}
概念
\exercisecheck{题目}{可靠性、可用性和可维护性各自关注什么？举例说明三者可能的取舍。}
\end{chapterexercise}

\begin{chapterexercise}{2}
错误分级
\exercisecheck{题目}{为什么“ECC 已纠正”仍需要记录地址和累计次数？}
\end{chapterexercise}

\begin{chapterexercise}{3}
故障传播
\exercisecheck{题目}{从内存控制器到可更换 DIMM，列出一次重复可纠正错误所需的观测字段。}
\end{chapterexercise}

\begin{chapterexercise}{4}
复位顺序
\exercisecheck{题目}{为一个仍可能 DMA 的设备设计局部复位顺序，并说明何时才能撤销 IOMMU 映射。}
\end{chapterexercise}

\begin{chapterexercise}{5}
迟到完成
\exercisecheck{题目}{设备复位后旧完成写回队列，如何避免它被识别为新请求的完成？}
\end{chapterexercise}

\begin{chapterexercise}{6}
验证
\exercisecheck{题目}{设计一次丢中断故障注入，说明预期检测、恢复和验收条件。}
\end{chapterexercise}

\begin{chapterexercise}{SECDED 分类}
分别判断以下组合：syndrome=0/overall=1，syndrome=5/overall=1，syndrome=5/overall=0，并给出数据动作。
\exercisecheck{检查要点}{双 bit 错误不能按 syndrome 误纠正}
\end{chapterexercise}

\begin{chapterexercise}{恢复验收}
设备 FLR 后重新枚举成功。列出证明它可以安全恢复 DMA 队列的五项不变量。
\exercisecheck{检查要点}{检查旧事务、地址映射、队列世代、中断与数据完整性}
\end{chapterexercise}

\section*{参考解答与要点}

\begin{chapteranswer}{1}
概念
\answercheck{解答要点}{可靠性关注结果正确，可用性关注服务连续，维护性关注定位与修复。发现风险后停机可提高正确性保障，却暂时降低可用性。}
\end{chapteranswer}

\begin{chapteranswer}{2}
错误分级
\answercheck{解答要点}{单次纠正隐藏了器件劣化趋势；重复集中在同一页、bank 或器件时，需要主动迁移、下线和维护。}
\end{chapteranswer}

\begin{chapteranswer}{3}
故障传播
\answercheck{解答要点}{至少含时间、物理地址、通道、DIMM、rank/bank、syndrome、纠正等级、首次状态、累计数和溢出标志。}
\end{chapteranswer}

\begin{chapteranswer}{4}
复位顺序
\answercheck{解答要点}{停止提交，屏蔽中断，等待或取消事务，确认设备停止 DMA，保存错误状态，再撤映射并复位；重建队列与映射后才恢复流量。}
\end{chapteranswer}

\begin{chapteranswer}{5}
迟到完成
\answercheck{解答要点}{在请求和完成中加入 generation/phase，复位时递增；清理旧队列并拒绝 generation 不匹配的完成。}
\end{chapteranswer}

\begin{chapteranswer}{6}
验证
\answercheck{解答要点}{屏蔽一次完成中断但保留完成状态；预期 watchdog/轮询发现队列无进展，驱动读取完成并恢复中断；验收包括无重复完成、无请求遗失和错误计数正确。}
\end{chapteranswer}

\begin{chapteranswer}{SECDED 分类}
0/1 表示额外总奇偶位错误，修正校验位；5/1 表示第 5 位单 bit 错误，翻转后交付并记录；5/0 表示检测到双 bit 错误，不得按位置 5 纠正。
\answercheck{必须保留的检查点}{syndrome 必须与 overall parity 联合解释}
\end{chapteranswer}

\begin{chapteranswer}{恢复验收}
确认旧 DMA 已静止，IOMMU 只含新映射，新 generation 拒绝旧完成，MSI-X 向量与新队列一致，描述符所有权与收发序号连续且无重复；随后再恢复上层流量。
\answercheck{必须保留的检查点}{重新枚举只证明配置访问恢复}
\end{chapteranswer}
